\section{本章小结}
\label{sec:rd-summary}

本章回答了开头的两个问题。下雨概率为$0.3$的记录，允许错$10\%$时每天只需$R(0.1)=h_2(0.3)-h_2(0.1)\approx0.412$\,bit，不到无损所需的一半。标准高斯测量值每个$1$\,bit时，长块编码的均方误差极限是$0.25$，逐个量化的最好结果是$0.363$。两处都是“允许失真的最低渐近码率”与“有限方案实际做到的”之间的比较。

这个比较涉及三个层次的对象，它们回答不同的问题，使用的标准必须相同：同一个源分布、同一个失真函数、期望失真、以每个源符号的bit数计的码率。
\begin{itemize}
  \item 有限方案，如第~\ref{chap:quantization}章的量化器或第~\ref{sec:rd-block-model}节的小块长码，可以直接执行，性能可以精确计算，但只是可达区域中的一个点。
  \item 信息率失真函数$R(D)$是单个符号上的凸优化：在平均失真不超过$D$的测试信道中最小化互信息。它非增、凸，在$[D_{\min},D_{\max}]$上严格递减，可以解析或数值地求出。
  \item 操作率失真函数$R_{\mathrm{op}}(D)$是所有块长、所有码本的极限。编码定理说明，对有限字母表、无记忆信源与有界可分失真，它与$R(D)$相等。
\end{itemize}

编码定理的两半依赖不同的工具。下界从“编号只有$M$个值”出发，经数据处理、独立性与链式法则拆成单个符号，再用互信息对测试信道的凸性合并；无记忆性与可分失真在其中缺一不可。可达性用随机码本：按最优测试信道的重构边缘撒点，大数律使信息密度集中在$R(D)$附近，码字数略多于$2^{nR(D)}$就能以高概率覆盖源块，再由平均推出好码本存在。

\begin{BookInlineTable}
  \centering
  \caption{贯穿例子中几类方案的比较。上半部分为$\operatorname{Bern}(0.3)$信源与Hamming失真，下半部分为标准高斯信源与平方失真。}
  \label{tab:rd-summary}
  \begin{tabular}{lcc}
    \toprule
    方案 & 码率（bit/符号） & 平均失真 \\
    \midrule
    逐位照传 & $1$ & $0$ \\
    块长$3$，码本$\{000,100,010,001\}$ & $0.667$ & $0.081$ \\
    块长$4$，最优$4$码字码本 & $0.5$ & $0.136$ \\
    率失真极限$R(D)$ & $0.412$ & $0.1$ \\
    率失真极限$D(R)$ & $0.5$ & $0.074$ \\
    \midrule
    两点Lloyd–Max量化器 & $1$ & $0.363$ \\
    均匀量化加索引熵编码（步长$1$） & $2.105$ & $0.077$ \\
    率失真极限$D(R)=2^{-2R}$ & $1$ & $0.25$ \\
    \bottomrule
  \end{tabular}
\end{BookInlineTable}

解析解都来自同一个两步法：一个对所有测试信道成立的信息下界，加一个达到它的反向测试信道。二元信源的下界来自错误指示的熵，高斯信源的下界来自条件高斯参考与KL非负，后者与固定方差的熵上界是同一个论证。独立高斯分量的码率按反注水分配，方差低于水平的分量不描述；相关分量先按协方差的特征方向旋转，再对特征值反注水。

有限概率表用Blahut–Arimoto算法计算。乘子$s$把失真约束变成罚项，最小点是曲线上斜率为$-s$的切点；交替地用“当前边缘乘以失真的指数权重再归一化”更新测试信道、用输出边缘更新$r$，目标单调下降，误差界$\max\log_2c$可以随时计算，目标误差至少按$1/T$缩小。

求出$R(D)$还没有给出编码器。它给出的是标准：一个有限方案离极限还有多远，例如高分辨率下均匀量化加熵编码约差$0.254$\,bit，这一段只能靠联合处理多个值去缩小；以及若用随机码本，码字应按什么分布生成。有限块长的最优码与低复杂度的搜索不由本章的定理给出。

描述一个信源需要$R(D)$的码率，这些bit还要经过有噪声的信道传出去。下一部分研究信道每次使用能可靠传送多少bit，即信道容量，再把所需的描述资源与可用的传输资源放在一起比较。
